% Variante de non-regression : diapositives d'une annexe (\slideappendix).
% Lue par check-appendix.sh : les lignes APPCHECK doivent se rendre avec des
% numeros en lettre (B.1, B.1.1), et le bandeau doit porter « Annexe B ». La
% seconde section verifie que les exemples repartent a 1 (B.2.1).
\documentclass[9pt,t]{beamer}
\usepackage[lang=fr, theme=ocots, math=analysis,
            institution={n7}, author={Prénom Nom}]{ocots}
\lstset{language=[LaTeX]TeX}

\title{Le template \texttt{ocots}}
\author{Prénom \textsc{Nom}}
\date{2025--2026}

\newcommand{\ocotsvariantnote}[1]{\begin{quote}\itshape Variante : #1\end{quote}}

\begin{document}

\begin{frame}[fragile]{Variante~: diapositives d'une annexe}
\begin{lstlisting}
\slideappendix{2}{Titre de l'annexe}
\end{lstlisting}
\ocotsvariantnote{\texttt{\textbackslash slideappendix} remplace
\texttt{\textbackslash slidechapter} pour une annexe du polycopié : le rang
\texttt{2} s'affiche en lettre (B) dans le bandeau, les sections et les
numéros des boîtes, qui reprennent ainsi ceux du polycopié.}
\end{frame}

\slideappendix{2}{Titre de l'annexe}
\slidetitlepage

\section{Dérivée d'une intégrale}

\begin{slide}{Numérotation d'une annexe}

\begin{theorem}[title={de dérivation de Lebesgue}, label=thm:app]
    Soit $f$ intégrable sur $\intervalcc{a}{b}$. Alors $x \mapsto \int_a^x f$ est dérivable presque partout, de
    dérivée $f$.
\end{theorem}

\begin{example}[label=exa:app]
    $f = \mathbb{1}_{\intervalcc{0}{1}}$ sur $\intervalcc{0}{2}$.
\end{example}

APPCHECK Section = \thesection

APPCHECK Théorème = Théorème~\ref{thm:app}

APPCHECK Exemple = Exemple~\ref{exa:app}

\end{slide}

\section{Fonctions absolument continues}

\begin{slide}{Remise à zéro par section}

\begin{example}[label=exa:app2]
    L'escalier de Cantor.
\end{example}

APPCHECK Exemple2 = Exemple~\ref{exa:app2}

\end{slide}

\begin{frame}{Sommaire}
\tableofcontents
\end{frame}

\end{document}
